\label{sec:performance-solver-behavior}

We now report numerical results for quadrature sensitivity, nonlinear stopping,
controlled globalization, and fixed-state tangent-action agreement. These experiments
clarify how the verified derivatives behave under the numerical choices used by
the solver. The available local studies support correctness and bounded
outcome statements, but not timing, scaling, or route-selection claims.

\subsection{Fixed-state quadrature sensitivity}
\label{subsec:quadrature-sensitivity}

We evaluate one prescribed coefficient vector and one deterministic unit
direction for each polynomial degree. The degree-associated tetrahedral rules
use 1, 11, and 24 points for $P_1$, $P_2$, and $P_4$, respectively; every
result is compared with the same 125-point Duffy rule. Let
$\Pi_q=E_q^{\mathrm{int}}-W_q^{\mathrm{ext}}$ and define
\begin{equation}
  s_{\Pi,125}:=\max\bigl(
  |E_{125}^{\mathrm{int}}|,|W_{125}^{\mathrm{ext}}|,|\Pi_{125}|,
  \operatorname{tiny}_{64}\bigr),
  \qquad
  \delta_\Pi(q):=\frac{|\Pi_q-\Pi_{125}|}{s_{\Pi,125}},
  \label{eq:quadrature-energy-scale}
\end{equation}
where $\operatorname{tiny}_{64}$ is the smallest positive normal binary64
number. We report, in order, $\delta_\Pi$, the relative free-residual defect
from \eqref{eq:verification-relative-defects}, the corresponding relative free
Hessian-action defect, and the $L^1$ difference of the absolute-quadrature-
weight branch fractions.

For $P_1$, these four discrepancies are $3.51\times10^{-16}$,
$2.79\times10^{-16}$, $3.60\times10^{-16}$, and
$7.11\times10^{-17}$. The affine-element integrands therefore agree to
roundoff. For $P_2$, the corresponding values are $7.11\times10^{-5}$,
$3.43\times10^{-3}$, $3.24\times10^{-2}$, and
$1.21\times10^{-2}$. For $P_4$, they increase to
$6.28\times10^{-4}$, $2.68\times10^{-2}$,
$6.91\times10^{-2}$, and $8.51\times10^{-3}$.

The energy discrepancy alone thus understates the residual and tangent
sensitivity for both higher-order spaces. Because the states are prescribed,
these comparisons do not measure endpoint stationarity or select a quadrature
policy. The aggregate branch-fraction differences also do not locate a branch
switch or establish differentiability there. A convergence conclusion requires
separately converged nonlinear endpoints under each rule.

\subsection{Mesh-aware stopping results}
\label{subsec:stopping-results}

Table~\ref{tab:stopping-local-status} summarizes the completed local stopping
matrix. All 45 locally feasible executions have checksum-identified outputs and
independent numerical rechecks. The endpoint gate admits 43 rows: all
Ginzburg--Landau,
hyperelasticity, and Plasticity3D nonlinear endpoints pass, while 13 of the 15
Plasticity3D fixed-state endpoints pass. Applying the stricter comparison
criteria admits 28 of the 45 local rows.

\begin{table}[t]
\centering
\caption{Fail-closed status of the local stopping-calibration tranche.}
\label{tab:stopping-local-status}
\begin{tabular}{lrrr}
\toprule
Family & Receipts & Admitted endpoints & Accepted comparisons \\
\midrule
Ginzburg--Landau endpoints & 8/8 & 8/8 & 5/8 \\
Hyperelasticity metric checks & 6/6 & 6/6 & 6/6 \\
Hyperelasticity nonlinear endpoints & 8/8 & 8/8 & 6/8 \\
Plasticity3D fixed-state solves & 15/15 & 13/15 & 8/15 \\
Plasticity3D nonlinear endpoints & 8/8 & 8/8 & 3/8 \\
\bottomrule
\end{tabular}
\begin{minipage}{0.96\linewidth}\small
All 45 serial rows are hash-bound and independently rechecked. Four P4 nonlinear rows and three MPI-consistency rows remain cluster-deferred. Closing them is necessary but not sufficient for a complete EXP-STOP-001 pass: separately hash-bound EXP-DISC evidence and final adjudication are also required. Timing, scaling, performance, and population claims are excluded.
\end{minipage}
\end{table}


Seven prescribed tests remain outside the local matrix: four $P_4$ nonlinear
calibrations and three publication-rank MPI-consistency comparisons. Their
absence prevents a complete mesh- and rank-aware stopping conclusion. The
local counts therefore verify the stated endpoint checks on the tested rows;
they do not identify a preferred stopping policy or support a timing claim.

\subsection{Controlled globalization outcomes}
\label{subsec:globalization-results}

The controlled design uses three prescribed starts and five separate
executions for each method and problem, giving the 60 two-rank
outcomes in Table~\ref{tab:globalization-local-status}. For the
Ginzburg--Landau problem on $L_5$, Newton--Armijo fails in all 15 executions, whereas
reduced trust--Armijo completes all 15. For the first hyperelastic load step on
$L_2$, both methods complete all 15 executions.

\begin{table}[t]
  \caption{Observed terminal outcomes for the frozen local controlled-globalization design.}
  \label{tab:globalization-local-status}
  \centering
  \begin{tabularx}{\linewidth}{C{1.55}C{1.15}C{0.65}C{0.65}}
    \toprule
    Problem & Method & Completed & Failed \\
    \midrule
    Ginzburg--Landau, $L_5$, 2 ranks & Newton--Armijo & 0/15 & 15/15 \\
    Ginzburg--Landau, $L_5$, 2 ranks & Reduced trust--Armijo & 15/15 & 0/15 \\
    Hyperelasticity, $L_2$, step 1, 2 ranks & Newton--Armijo & 15/15 & 0/15 \\
    Hyperelasticity, $L_2$, step 1, 2 ranks & Reduced trust--Armijo & 15/15 & 0/15 \\
    \bottomrule
  \end{tabularx}
  \begin{minipage}{0.96\linewidth}\small
    Counts are outcomes for three prescribed starts and five process repetitions per method. Timing and performance ordering are excluded. The starts are deterministic sensitivity instances, so no robustness generalization is made. A paired method comparison is admissible only when both methods terminate at the same independently identified endpoint.
  \end{minipage}
\end{table}


The common-start identity check passes, but the paired endpoint-identity gate
does not. The Ginzburg--Landau pairs are incomplete because Newton--Armijo does
not terminate successfully, and the saved hyperelastic endpoints from the two
methods have different identities. Consequently, the experiment supports only
the displayed completion and failure counts for these deterministic starts. It
does not support a paired method ranking, a robustness probability, or a timing
comparison.

\subsection{Fixed-state tangent-action agreement}
\label{subsec:route-map-results}

The checksum-verified workstation study compares element automatic differentiation,
constitutive automatic differentiation, and colored sparse recovery for
$P_1(L_1)$ and $P_2(L_1)$ at elastic and mixed states. The 12 balanced blocks
contain 36 sequential route executions, and each of the six route orders occurs
twice. All blocks complete, and the declared output inventory passes its
independent checksum verification.

Across four deterministic probes per route and state, the largest relative
tangent-action discrepancy is $2.34\times10^{-16}$. The saved states and the
common element-AD residuals agree exactly at stored precision. This finite study
therefore corroborates four-probe tangent-action agreement on the tested serial
states. Since it contains neither a distributed timing design nor the full
degree--rank coverage, it defines no finite performance map, crossover
predictor, or route ranking.

Together, these experiments quantify higher-order quadrature sensitivity,
establish the local extent of the stopping checks, record bounded globalization
outcomes, and verify serial four-probe tangent-action agreement. Their explicit
scope excludes
nonlinear quadrature convergence, a general stopping or globalization policy,
parallel scaling, and performance-based route selection.

\FloatBarrier
